\section{Liquor Composition Along a Multicomponent Eutectic Cascade}

In a multicomponent eutectic purification cascade, the system evolves through a
sequence of nodes. At each node $k$, the eutectic solver returns two quantities:

\begin{itemize}
    \item The eutectic solid composition:
    

\[
        X^{(k)}_{\mathrm{eu}} = \{ X^{(k)}_{\mathrm{eu},i} \}_{i=1}^N
    \]


    \item The remaining mixture after subtracting the eutectic:
    

\[
        X^{(k)}_{\mathrm{mix}} = \{ X^{(k)}_{\mathrm{mix},i} \}_{i=1}^N
    \]


\end{itemize}

The remaining mixture $X^{(k)}_{\mathrm{mix}}$ is obtained by subtracting the
eutectic mass from the current mixture and renormalising. This cumulative
renormalisation implicitly contains the full history of all previous eutectic
removals. Therefore, $X^{(k)}_{\mathrm{mix}}$ represents the true liquor
composition at node $k$.

\subsection{Canonical Segments}

The canonical cascade is constructed by pairing successive nodes into segments.
For segment $s$ spanning nodes $k$ and $k+1$, we define:



\[
    X^{(s)}_{\mathrm{liq,start}} = X^{(k)}_{\mathrm{mix}}, \qquad
    X^{(s)}_{\mathrm{liq,end}}   = X^{(k+1)}_{\mathrm{mix}}.
\]



These two compositions are physically valid equilibrium liquors produced by the
solver. They lie in the simplex, satisfy mass balance, and require no further
derivation.

The solid composition associated with segment $s$ is the eutectic at the
\emph{end} node:



\[
    X^{(s)}_{\mathrm{sol}} = X^{(k+1)}_{\mathrm{eu}}.
\]



This is consistent with the gradient/intercept construction, where the solid
composition is determined by the eutectic geometry at the segment boundary.

\subsection{Liquor Composition at Arbitrary Recovery}

For any recovery value $R$ lying within segment $s$, the liquor composition is
obtained by linear interpolation between the segment endpoints:



\[
    \lambda = \frac{R - R_{\mathrm{start}}^{(s)}}
                   {R_{\mathrm{end}}^{(s)} - R_{\mathrm{start}}^{(s)}},
\]





\[
    X^{(R)}_{\mathrm{liq},i}
    = (1 - \lambda)\, X^{(s)}_{\mathrm{liq,start},i}
      + \lambda\, X^{(s)}_{\mathrm{liq,end},i}.
\]



Because both endpoint compositions are valid equilibrium liquors, the
interpolated composition $X^{(R)}_{\mathrm{liq}}$ remains within the feasible
composition simplex and satisfies:



\[
    X^{(R)}_{\mathrm{liq},i} \ge 0, \qquad
    \sum_{i=1}^N X^{(R)}_{\mathrm{liq},i} = 1.
\]



\subsection{Why the Lever Rule is Not Used}

In multicomponent, non-ideal systems, the liquor composition cannot be derived
from the solid composition via a binary lever rule:



\[
    X_{\mathrm{liq},i}
    = \frac{X_{c,i} - R\,X_{\mathrm{sol},i}}{1 - R}.
\]



This expression does not account for the cumulative history of eutectic
subtractions and fails to represent the true equilibrium liquor in systems with
more than two components or non-ideal interactions. For this reason, the liquor
composition is taken directly from the solver outputs $X^{(k)}_{\mathrm{mix}}$
and interpolated along canonical segments.

\subsection{Summary}

\begin{itemize}
    \item The solid compositions $X^{(s)}_{\mathrm{sol}}$ come directly from the
          eutectic solver and correspond to the gradient/intercept geometry at
          each node.
    \item The liquor compositions $X^{(s)}_{\mathrm{liq,start}}$ and
          $X^{(s)}_{\mathrm{liq,end}}$ are the renormalised remaining mixtures
          after each eutectic removal.
    \item Slice liquor compositions are obtained by interpolation between these
          physically valid endpoint liquors.
\end{itemize}

This construction ensures that all slice compositions are physically meaningful,
mass-balanced, and consistent with the underlying thermodynamic model.
